\documentclass[]{article}
\usepackage[margin=1in]{geometry}
\usepackage{amsmath,amssymb,amsthm}
\usepackage[most]{tcolorbox}
\usepackage{xcolor}
\usepackage{enumitem}
\usepackage{fancyhdr}

% Page setup
\pagestyle{fancy}
\fancyhf{}
\rhead{AMC12/COMC Problem Analysis}
\lhead{\today}

% Color definitions
\definecolor{problemconfig}{RGB}{230, 242, 255}
\definecolor{strategic}{RGB}{255, 230, 242}
\definecolor{tactical}{RGB}{242, 255, 230}
\definecolor{techniques}{RGB}{255, 242, 230}
\definecolor{verification}{RGB}{255, 230, 230}
\definecolor{solution}{RGB}{230, 255, 255}
\definecolor{competition}{RGB}{242, 242, 255}
\definecolor{gold}{RGB}{218, 165, 32}

% Box styles
\newtcolorbox{problemconfigbox}{
    colback=problemconfig,
    colframe=blue,
    title=Problem Configuration Analysis,
    fonttitle=\bfseries,
    arc=3pt,
    boxrule=1pt,
    breakable
}

\newtcolorbox{strategicbox}{
    colback=strategic,
    colframe=purple,
    title=Strategic Framework Application,
    fonttitle=\bfseries,
    arc=3pt,
    boxrule=1pt,
    breakable
}

\newtcolorbox{tacticalbox}{
    colback=tactical,
    colframe=green,
    title=Tactical Methodology Selection,
    fonttitle=\bfseries,
    arc=3pt,
    boxrule=1pt,
    breakable
}

\newtcolorbox{techniquesbox}{
    colback=techniques,
    colframe=orange,
    title=Conversion Techniques Application,
    fonttitle=\bfseries,
    arc=3pt,
    boxrule=1pt,
    breakable
}

\newtcolorbox{verificationbox}{
    colback=verification,
    colframe=red,
    title=Verification Strategy Design,
    fonttitle=\bfseries,
    arc=3pt,
    boxrule=1pt,
    breakable
}

\newtcolorbox{solutionbox}{
    colback=solution,
    colframe=cyan,
    title=Solution Roadmap,
    fonttitle=\bfseries,
    arc=3pt,
    boxrule=1pt,
    breakable
}

\newtcolorbox{competitionbox}{
    colback=competition,
    colframe=purple,
    title=Competition Strategy Integration,
    fonttitle=\bfseries,
    arc=3pt,
    boxrule=1pt,
    breakable
}

\newtcolorbox{keyinsightbox}{
    colback=yellow!20,
    colframe=gold,
    title=Key Insight,
    fonttitle=\bfseries,
    arc=3pt,
    boxrule=2pt,
    leftrule=5pt,
    breakable
}

\newtcolorbox{warningbox}{
    colback=red!10,
    colframe=red!50,
    title=Warning: Common Pitfall,
    fonttitle=\bfseries,
    arc=3pt,
    boxrule=2pt,
    leftrule=5pt,
    breakable
}

\newtcolorbox{tipbox}{
    colback=green!10,
    colframe=green!50,
    title=Pro Tip,
    fonttitle=\bfseries,
    arc=3pt,
    boxrule=2pt,
    leftrule=5pt,
    breakable
}

\begin{document}

\title{AMC12B 2017 Problem 14: Strategic Analysis\\
\large Ice Cream Cone and Frustum Volume}
\author{Competition Math Framework}
\date{\today}
\maketitle

\section*{Problem Statement}
An ice-cream novelty item consists of a cup in the shape of a 4-inch-tall frustum of a right circular cone, with a 2-inch-diameter base at the bottom and a 4-inch-diameter base at the top, packed solid with ice cream, together with a solid cone of ice cream of height 4 inches, whose base, at the bottom, is the top base of the frustum. What is the total volume of the ice cream, in cubic inches?

$$\textbf{(A)}\ 8\pi \qquad \textbf{(B)}\ \frac{28\pi}{3} \qquad \textbf{(C)}\ 12\pi \qquad \textbf{(D)}\ 14\pi \qquad \textbf{(E)}\ \frac{44\pi}{3}$$

\begin{problemconfigbox}
\subsection*{A. Problem Classification}
\begin{itemize}[leftmargin=*]
    \item \textbf{Problem Type}: To Find (compute exact numerical volume)
    \item \textbf{Difficulty Band}: Mid-band (Problem 14 of 25)
    \item \textbf{Competition Context}: AMC12B 2017
    \item \textbf{Mathematical Domain}: Solid Geometry (Volume Calculations)
\end{itemize}

\subsection*{B. Problem Structure Identification}
\begin{itemize}[leftmargin=*]
    \item \textbf{Given Information}: 
    \begin{itemize}
        \item Frustum: height $h_1 = 4$ inches, bottom diameter $d_1 = 2$ inches (radius $r_1 = 1$), top diameter $d_2 = 4$ inches (radius $r_2 = 2$)
        \item Cone on top: height $h_2 = 4$ inches, base diameter $d_3 = 4$ inches (radius $r_3 = 2$)
        \item Cone's base sits on frustum's top
    \end{itemize}
    \item \textbf{Constraints}: 
    \begin{itemize}
        \item All shapes are right circular (axis perpendicular to bases)
        \item Frustum is part of a cone (linear taper from bottom to top)
        \item Answer must be in terms of $\pi$
    \end{itemize}
    \item \textbf{Objective}: Find total volume $V_{\text{total}} = V_{\text{frustum}} + V_{\text{cone}}$
    \item \textbf{Hidden Assumptions}: 
    \begin{itemize}
        \item Frustum formula not explicitly given (must recall or derive)
        \item Similar triangles relationship in frustum structure
    \end{itemize}
    \item \textbf{Answer Choice Leverage}: 
    \begin{itemize}
        \item All choices involve $\pi$ and are relatively simple expressions
        \item Choices (B) and (E) have denominators of 3 (suggesting $\frac{1}{3}$ from cone formula)
        \item Choice (E) is largest, suggesting additive answer
        \item Can verify by checking if answer is reasonable for given dimensions
    \end{itemize}
\end{itemize}

\subsection*{C. Strategic Orientation}
\begin{itemize}[leftmargin=*]
    \item \textbf{Hypothesis}: Frustum can be computed as difference of two cones
    \item \textbf{Conclusion}: Total volume equals frustum volume plus top cone volume
    \item \textbf{Notation}: 
    \begin{itemize}
        \item $V_f$ = volume of frustum
        \item $V_c$ = volume of top cone  
        \item $r_1, r_2$ = radii at bottom and top of frustum
        \item $h_1, h_2$ = heights of frustum and cone
    \end{itemize}
    \item \textbf{Intuitive Approach}: 
    \begin{itemize}
        \item Cone formula is straightforward: $V = \frac{1}{3}\pi r^2 h$
        \item Frustum requires either: (1) direct formula, or (2) difference of cones method
        \item Difference method: imagine completing the frustum to a full cone
    \end{itemize}
    \item \textbf{Cross-Domain Potential}: Could use calculus integration with $\pi \int_0^4 r(h)^2 \, dh$, but algebraic approach is simpler for AMC12
\end{itemize}
\end{problemconfigbox}

\begin{strategicbox}
\subsection*{A. Psychological Strategies}
\begin{itemize}[leftmargin=*]
    \item \textbf{Mental Toughness}: This is a mid-band problem (14/25), so expect 4-6 minutes of work. If frustum formula isn't recalled, derive it systematically rather than giving up
    \item \textbf{Creativity Cultivation}: Multiple approaches exist (direct formula vs. cone difference). If one stalls, switch methods
    \item \textbf{Iterative Refinement}: After computing frustum volume, double-check similar triangles setup before proceeding to final addition
\end{itemize}

\subsection*{B. Investigation Strategies}
\begin{itemize}[leftmargin=*]
    \item \textbf{Startup Orientation}: Recognize this as volume decomposition problem. Sketch the shape to visualize frustum + cone structure
    \item \textbf{Get Your Hands Dirty}: Start with what's known: top cone has clear dimensions ($r=2$, $h=4$)
    \item \textbf{Penultimate Step}: Final step will be addition $V_f + V_c$, so focus on getting frustum volume correct
    \item \textbf{Wishful Thinking}: ``I wish I remembered the frustum formula'' $\rightarrow$ derive it via cone subtraction
    \item \textbf{Make It Easier}: Top cone is easier than frustum, so compute it first for a small confidence boost
    \item \textbf{Conjecture via Small Cases}: Verify formula with simple example: if frustum has $r_1=r_2$, it becomes a cylinder ($V=\pi r^2 h$)
    \item \textbf{Visualization}: Draw cross-section showing frustum as part of cone with vertex below bottom base
\end{itemize}

\subsection*{C. Late-Band Strategic Playbook}
Not directly applicable (this is mid-band P14), but general principles:
\begin{itemize}[leftmargin=*]
    \item \textbf{Answer-Choice Leverage}: Check if answer choices help eliminate options (e.g., too small/large)
    \item \textbf{Dimensional Analysis}: Answer must have units of $\text{inches}^3$ and factor of $\pi$
\end{itemize}

\subsection*{D. Argument Construction Strategy}
\begin{itemize}[leftmargin=*]
    \item \textbf{Direct Proof}: Compute each volume component and add
    \item \textbf{Geometric Construction}: Extend frustum lines to find apex of complete cone
    \item \textbf{Similar Triangles}: Use proportionality to find height of small cone above frustum
\end{itemize}

\subsection*{E. Cross-Domain Translation}
\begin{itemize}[leftmargin=*]
    \item \textbf{Draw a Picture}: Essential for understanding frustum as part of larger cone
    \item \textbf{Algebraic vs. Geometric}: While calculus integration is possible, geometric decomposition is cleaner
    \item \textbf{Coordinate System}: Could place origin at frustum base, but not necessary for this problem
\end{itemize}
\end{strategicbox}

\begin{tacticalbox}
\subsection*{A. Core Mathematical Tactics}
\begin{itemize}[leftmargin=*]
    \item \textbf{Primary Tactic}: Volume Decomposition (TM31-VolumeDecompose)
    \item \textbf{Supporting Tactics}:
    \begin{itemize}
        \item Similar Triangles (T3-SimilarTriangles) for finding cone apex
        \item Algebraic Manipulation for combining terms
    \end{itemize}
\end{itemize}

\subsection*{B. Domain-Specific Tactics}
\begin{itemize}[leftmargin=*]
    \item \textbf{Solid Geometry Fundamentals}:
    \begin{itemize}
        \item Cone volume: $V = \frac{1}{3}\pi r^2 h$
        \item Frustum as difference of cones: $V_{\text{frustum}} = V_{\text{large}} - V_{\text{small}}$
        \item Direct frustum formula: $V = \frac{\pi h}{3}(r_1^2 + r_1 r_2 + r_2^2)$
    \end{itemize}
    \item \textbf{Similar Triangles}:
    \begin{itemize}
        \item In cross-section, frustum sides form similar triangles with full cone
        \item If frustum has bottom radius $r_1$, top radius $r_2$, height $h$
        \item Then small cone removed has height ratio $\frac{r_1}{r_2-r_1} \cdot h$
    \end{itemize}
\end{itemize}

\subsection*{C. Crossover Techniques}
\begin{itemize}[leftmargin=*]
    \item \textbf{Calculus Alternative}: $V_{\text{frustum}} = \pi \int_0^4 \left(\frac{h}{4}+1\right)^2 dh$
    \item \textbf{Coordinate Geometry}: Place cone apex at origin and integrate, but overkill for AMC12
\end{itemize}
\end{tacticalbox}

\begin{techniquesbox}
\subsection*{A. Recognition Index}
\begin{itemize}[leftmargin=*]
    \item \textbf{Pattern Recognition}: Composite solid with two components
    \item \textbf{Formula Recognition}: Cone volume formula is standard; frustum requires recall or derivation
    \item \textbf{Structure Recognition}: Frustum is truncated cone (part of larger cone)
\end{itemize}

\subsection*{B. Technique Selection}
\begin{itemize}[leftmargin=*]
    \item \textbf{Primary Technique}: TM31-VolumeDecompose (Volume Decomposition)
    \item \textbf{Confidence Level}: 95\% (standard technique for composite solids)
    \item \textbf{Alternative Techniques}:
    \begin{itemize}
        \item Direct frustum formula (requires memorization)
        \item Calculus integration (slower but verifiable)
    \end{itemize}
\end{itemize}

\subsection*{C. Technique Integration}
\begin{itemize}[leftmargin=*]
    \item \textbf{Step 1}: Compute top cone volume (straightforward)
    \item \textbf{Step 2}: Compute frustum volume via cone subtraction
    \begin{itemize}
        \item Find dimensions of large cone (extending frustum downward)
        \item Find dimensions of small cone (removed part above frustum)
        \item Subtract: $V_f = V_{\text{large}} - V_{\text{small}}$
    \end{itemize}
    \item \textbf{Step 3}: Add components: $V_{\text{total}} = V_f + V_c$
\end{itemize}

\subsection*{D. Conflict Resolution}
\begin{itemize}[leftmargin=*]
    \item \textbf{Potential Issue}: Frustum formula not recalled
    \item \textbf{Resolution}: Derive using similar triangles and cone subtraction
    \item \textbf{Verification}: Check that frustum formula gives same result as direct computation
\end{itemize}
\end{techniquesbox}

\begin{verificationbox}
\subsection*{A. Level Selection with Decision Tree}
\begin{itemize}[leftmargin=*]
    \item \textbf{Time Available}: 4-6 minutes for mid-band problem
    \item \textbf{Confidence}: High (standard formulas and clear structure)
    \item \textbf{Verification Level}: Level 3 (Sanity + Primary)
    \begin{itemize}
        \item Level 1: Quick check (units, sign)
        \item Level 2: Recompute final addition
        \item Level 3: Verify frustum computation with alternative method
    \end{itemize}
\end{itemize}

\subsection*{B. Specialized Verification Methods}
\begin{itemize}[leftmargin=*]
    \item \textbf{Dimensional Analysis}:
    \begin{itemize}
        \item All terms must have units of $\text{inches}^3$
        \item Factor of $\pi$ must appear (circular bases)
        \item Answer should be between $12\pi$ and $20\pi$ (rough estimate)
    \end{itemize}
    \item \textbf{Answer Choice Verification}:
    \begin{itemize}
        \item Top cone: $V_c = \frac{1}{3}\pi (2)^2 (4) = \frac{16\pi}{3} \approx 16.76$ cubic inches
        \item Frustum should be larger than cone (wider base)
        \item Choice (E) $\frac{44\pi}{3} \approx 46.08$ is reasonable
    \end{itemize}
    \item \textbf{Algebraic Verification}:
    \begin{itemize}
        \item Recompute frustum using direct formula: $V = \frac{\pi h}{3}(r_1^2 + r_1 r_2 + r_2^2)$
        \item $V_f = \frac{\pi \cdot 4}{3}(1^2 + 1 \cdot 2 + 2^2) = \frac{4\pi}{3}(1+2+4) = \frac{28\pi}{3}$
        \item Matches our cone-subtraction result
    \end{itemize}
\end{itemize}

\subsection*{C. Confidence Calibration}
\begin{itemize}[leftmargin=*]
    \item \textbf{Initial Confidence}: 85\% (familiar problem type)
    \item \textbf{Post-Calculation Confidence}: 95\% (clean arithmetic, verified formula)
    \item \textbf{Post-Verification Confidence}: 98\% (two methods agree, answer choice reasonable)
\end{itemize}
\end{verificationbox}

\begin{solutionbox}
\subsection*{A. Step-by-Step Solution Plan}

\textbf{Phase 1: Compute Top Cone Volume (1 minute)}

The cone sits on top with radius $r = 2$ inches and height $h = 4$ inches.

Using cone formula:
$$V_c = \frac{1}{3}\pi r^2 h = \frac{1}{3}\pi (2)^2 (4) = \frac{16\pi}{3}$$

\textbf{Checkpoint 1}: Verify $\frac{16\pi}{3}$ is reasonable for 4-inch tall cone with 2-inch radius.

\textbf{Phase 2: Compute Frustum Volume via Cone Subtraction (3 minutes)}

The frustum has:
\begin{itemize}
    \item Bottom radius: $r_1 = 1$ inch
    \item Top radius: $r_2 = 2$ inches
    \item Height: $h_1 = 4$ inches
\end{itemize}

\textbf{Step 2a}: Find height of complete large cone

Using similar triangles in cross-section: if we extend the frustum downward to apex, let total height be $H$.

From similar triangles:
$$\frac{r_1}{H-h_1} = \frac{r_2}{H}$$
$$\frac{1}{H-4} = \frac{2}{H}$$
$$H = 2(H-4)$$
$$H = 2H - 8$$
$$H = 8 \text{ inches}$$

\textbf{Step 2b}: Compute large cone volume

Large cone: radius $r_2 = 2$, height $H = 8$
$$V_{\text{large}} = \frac{1}{3}\pi (2)^2 (8) = \frac{32\pi}{3}$$

\textbf{Step 2c}: Compute small cone volume (removed part)

Small cone: radius $r_1 = 1$, height $H - h_1 = 4$
$$V_{\text{small}} = \frac{1}{3}\pi (1)^2 (4) = \frac{4\pi}{3}$$

\textbf{Step 2d}: Compute frustum volume

$$V_f = V_{\text{large}} - V_{\text{small}} = \frac{32\pi}{3} - \frac{4\pi}{3} = \frac{28\pi}{3}$$

\textbf{Checkpoint 2}: Note that $\frac{28\pi}{3}$ matches answer choice (B), but this is only the frustum!

\textbf{Phase 3: Compute Total Volume (1 minute)}

$$V_{\text{total}} = V_f + V_c = \frac{28\pi}{3} + \frac{16\pi}{3} = \frac{44\pi}{3}$$

\textbf{Checkpoint 3}: This matches answer choice (E).

\subsection*{B. Alternative Approaches}

\textbf{Alternative 1: Direct Frustum Formula}

Using $V = \frac{\pi h}{3}(r_1^2 + r_1 r_2 + r_2^2)$:

$$V_f = \frac{\pi \cdot 4}{3}(1^2 + 1 \cdot 2 + 2^2) = \frac{4\pi}{3}(1 + 2 + 4) = \frac{4\pi}{3} \cdot 7 = \frac{28\pi}{3}$$

Same result! Then add cone volume: $\frac{28\pi}{3} + \frac{16\pi}{3} = \frac{44\pi}{3}$

\textbf{Alternative 2: Calculus Integration}

For frustum, radius varies linearly: $r(h) = 1 + \frac{h}{4}$ for $0 \leq h \leq 4$

$$V_f = \pi \int_0^4 r(h)^2 \, dh = \pi \int_0^4 \left(1 + \frac{h}{4}\right)^2 dh$$

$$= \pi \int_0^4 \left(1 + \frac{h}{2} + \frac{h^2}{16}\right) dh$$

$$= \pi \left[h + \frac{h^2}{4} + \frac{h^3}{48}\right]_0^4$$

$$= \pi \left(4 + 4 + \frac{64}{48}\right) = \pi \left(8 + \frac{4}{3}\right) = \frac{28\pi}{3}$$

Again matches!

\subsection*{C. Critical Pitfalls and How to Avoid Them}

\textbf{Pitfall 1}: Forgetting to add the top cone volume
\begin{itemize}
    \item \textbf{Warning}: The frustum volume $\frac{28\pi}{3}$ appears as answer choice (B) --- this is a trap!
    \item \textbf{Prevention}: Re-read problem to confirm both frustum AND cone are included
    \item \textbf{Recovery}: After selecting answer, verify by re-reading problem statement
\end{itemize}

\textbf{Pitfall 2}: Using diameter instead of radius
\begin{itemize}
    \item \textbf{Warning}: Problem gives diameters (2 and 4 inches), but formulas use radius
    \item \textbf{Prevention}: Immediately convert: $r_1 = 1$, $r_2 = 2$, $r_3 = 2$
    \item \textbf{Recovery}: If answer doesn't match choices, check radius vs. diameter
\end{itemize}

\textbf{Pitfall 3}: Incorrect similar triangles setup
\begin{itemize}
    \item \textbf{Warning}: Setting up ratio incorrectly when finding large cone height
    \item \textbf{Prevention}: Draw clear diagram with labels; verify $\frac{r_1}{H-h_1} = \frac{r_2}{H}$
    \item \textbf{Recovery}: Use alternative direct frustum formula as verification
\end{itemize}

\textbf{Pitfall 4}: Arithmetic errors in fraction addition
\begin{itemize}
    \item \textbf{Warning}: $\frac{28\pi}{3} + \frac{16\pi}{3}$ could be miscalculated
    \item \textbf{Prevention}: Factor out $\frac{\pi}{3}$ first: $\frac{\pi}{3}(28+16) = \frac{44\pi}{3}$
    \item \textbf{Recovery}: Recompute addition; verify against answer choices
\end{itemize}
\end{solutionbox}

\begin{competitionbox}
\subsection*{A. Time Management}
\begin{itemize}[leftmargin=*]
    \item \textbf{Target Time}: 4-6 minutes for Problem 14
    \item \textbf{Time Allocation}:
    \begin{itemize}
        \item Problem reading and setup: 0.5 min
        \item Top cone calculation: 1 min
        \item Frustum calculation: 2.5 min
        \item Addition and verification: 1 min
        \item Total: $\sim$5 minutes
    \end{itemize}
    \item \textbf{Actual Complexity}: Standard mid-band problem, should be comfortable within time
\end{itemize}

\subsection*{B. Problem Prioritization Strategy}
\begin{itemize}[leftmargin=*]
    \item \textbf{Accessibility}: High (clear problem statement, standard formulas)
    \item \textbf{Domain Comfort}: If comfortable with 3D geometry, tackle immediately
    \item \textbf{Point Value}: All AMC12 problems worth 6 points, so no special priority
    \item \textbf{Strategic Decision}: Good problem to solve early for confidence boost
\end{itemize}

\subsection*{C. Risk Management}
\begin{itemize}[leftmargin=*]
    \item \textbf{Confidence Threshold}: 95\% confidence after verification
    \item \textbf{Flag for Review}: No (high confidence, clean calculation)
    \item \textbf{Move-On Decision}: Proceed to next problem after brief check
\end{itemize}

\subsection*{D. Abandonment Criteria}
\textbf{Soft Abandon Triggers}:
\begin{itemize}[leftmargin=*]
    \item If frustum formula not recalled after 2 minutes $\rightarrow$ use similar triangles derivation
    \item If arithmetic becomes messy $\rightarrow$ slow down and recompute carefully
\end{itemize}

\textbf{Hard Abandon Triggers}:
\begin{itemize}[leftmargin=*]
    \item None applicable (this is straightforward mid-band problem)
    \item If spending $>8$ minutes, something is fundamentally wrong $\rightarrow$ skip and return
\end{itemize}

\textbf{Strategic Note}: This problem should not be abandoned. If stuck, use alternative method (direct frustum formula or calculus).
\end{competitionbox}

\begin{keyinsightbox}
\textbf{Key Insight}: The frustum is the difference between two cones with radii in ratio 2:1. Using similar triangles, the large cone has height 8 (twice the frustum height), making all calculations clean. The trap answer (B) $\frac{28\pi}{3}$ is just the frustum volume --- don't forget to add the top cone!
\end{keyinsightbox}

\begin{warningbox}
\textbf{Warning}: Answer choice (B) $\frac{28\pi}{3}$ is the frustum volume alone. This is a deliberate trap for students who forget to include the top cone. Always re-read the problem to ensure all components are included in the final answer.
\end{warningbox}

\begin{tipbox}
\textbf{Pro Tip}: For frustum problems, quickly check if the direct formula $V = \frac{\pi h}{3}(r_1^2 + r_1 r_2 + r_2^2)$ is faster than similar triangles. With simple integer radii like 1 and 2, direct computation is often cleaner. However, knowing both methods provides built-in verification!
\end{tipbox}

\section*{Final Answer}
$$\boxed{\textbf{(E)}\ \frac{44\pi}{3}}$$

\section*{Confidence Assessment}
\begin{itemize}[leftmargin=*]
    \item \textbf{Confidence Level}: 98\% (verified with multiple methods, answer reasonable)
    \item \textbf{Verification Applied}: Level 3 (Sanity + Primary + Alternative method check)
    \item \textbf{Flag for Review}: No (high confidence, clean arithmetic)
    \item \textbf{Time Spent}: 5 minutes (within target range for Problem 14)
\end{itemize}

\end{document}